% Bài làm buổi 20/10 — Gạo
\begin{problem}[9DS-PT2-003]
a) $x^2 - 6x + 4 = 0$ có $b' = -3$, $\Delta' = (-3)^2 - 1 \cdot 4 = 5 > 0$
Phương trình có hai nghiệm $x_1 = 3 + \sqrt{5}$; $x_2 = 3 - \sqrt{5}$
b) $9x^2 - 12x + 4 = 0$ có $b' = -6$, $\Delta' = (-6)^2 - 9 \cdot 4 = 0$
Phương trình có nghiệm kép $x_1 = x_2 = \dfrac{6}{9} = \dfrac{2}{3}$
c) $2x^2 + 3x + 5 = 0$ có $\Delta = 3^2 - 4 \cdot 2 \cdot 5 = 9 - 40 = -31 < 0$
Phương trình vô nghiệm.
\end{problem}

\begin{problem}[9DS-PT2-004]
a) $5x^2 - 7x + 2 = 0$ có $\Delta = (-7)^2 - 4 \cdot 5 \cdot 2 = 49 - 40 = 9$
$x_1 = \dfrac{7 + 3}{10} = 1$; $x_2 = \dfrac{7 - 3}{10} = \dfrac{2}{5}$
b) $2026x^2 + 2025x - 1 = 0$ có $a - b + c = 2026 - 2025 - 1 = 0$
Nên phương trình có nghiệm $x_1 = -1$; $x_2 = -\dfrac{c}{a} = \dfrac{1}{2026}$
\end{problem}

\begin{problem}[9DS-HSB2-002]
a) Đồ thị đi qua $A(-2; 2)$ nên $2 = a \cdot (-2)^2$
$4a = 2$, $a = \dfrac{1}{2}$
b) Với $a = \dfrac{1}{2}$ ta có $y = \dfrac{1}{2}x^2$.
Thay $x = 4$: $y = \dfrac{1}{2} \cdot 16 = 8$ nên $B$ thuộc đồ thị.
Thay $x = -1$: $y = \dfrac{1}{2} \cdot (-1)^2 = \dfrac{1}{2} \neq -\dfrac{1}{2}$ nên $C$ không thuộc đồ thị.
Thay $x = \sqrt{6}$ thì thấy $D$ thuộc đồ thị.
\end{problem}
